\section{Introduction}
\label{sec:intro}

Agents built on \acp{LLM} use long-term memory to retain information across
interactions, from personal assistance to simulated societies and
multi-episode tasks~\cite{park2023generative,packer2023memgpt,sumers2024coala,zhang2025survey}.
Their histories can exceed the model's context window; even within that
window, access to relevant information can depend on context length and
position~\cite{liu2024lost,hsieh2024ruler}. External memory makes these
histories available beyond a single interaction, but leaves a retrieval
problem: which evidence should the agent read to answer the current
question?

This problem is particularly demanding when an answer depends on facts
recorded in different sessions, connected through intermediate entities,
or associated with a particular period. Long-term conversational
benchmarks include precisely these forms of cross-session and temporal
reasoning~\cite{maharana2024locomo,wu2025longmemeval}. In such cases,
relevance is not solely a property of individual records. Two records
may jointly establish an answer even when one is only weakly related to
the question in isolation. Conversely, several highly relevant records
may repeat the same information while omitting a necessary link.

Ranking and iterative retrieval offer complementary responses to this
challenge. Ranking methods score memory records or propagate relevance
through a structured store before selecting a bounded
context~\cite{park2023generative,zhong2024memorybank,chhikara2025mem0,gutierrez2024hipporag}.
Their selection criteria need not explicitly test whether the retrieved
records satisfy the relations required by the answer. Iterative methods
use intermediate reasoning or model judgments to guide subsequent
retrieval~\cite{trivedi2023ircot,asai2024selfrag,jeong2024adaptiverag,yan2025gam}.
They can revise the search as evidence accumulates, but a model's
assessment of sufficiency does not itself identify a set of facts that
satisfies a specified relational condition. We investigate a complementary
approach in which these conditions are represented and executed explicitly.

Conjunctive queries provide a suitable formalism for this purpose.
A query specifies relations whose arguments must admit a consistent
assignment~\cite{chandra1977conjunctive}. A \emph{witness} is a set of
facts that supports one such assignment and can therefore explain a
query answer through its provenance~\cite{buneman2001why,green2007provenance}.
Applied to memory retrieval, this perspective separates the linguistic
task of specifying what to find from the computational task of finding
facts that jointly meet that specification. The resulting support can
be inspected independently of the reader's final response.

We introduce \sys, a retrieval framework organized into memory
construction and memory retrieval. During construction, the system
extracts self-contained propositions, records their relational
structure and temporal information, and retains links to the source
dialogue. During retrieval, an \ac{LLM} writes a conjunctive plan
conditioned on retrieved evidence. The plan specifies an answer type
and output operator, together with a temporal reference and ranking
weights. An executor searches for consistent bindings over the stored
graph without \ac{LLM} calls. A conditional verifier checks candidate
answers against source utterances, and execution gaps or rejections
guide a bounded number of planning cycles. Accepted witnesses then
guide the selection of dated propositions and passage summaries for
the reader.

The key idea is to make relational requirements part of retrieval
itself. A plan states which connections matter, and a witness identifies
the memories that realize those connections. This structure also gives
replanning a concrete basis: the next attempt can address a missing
relation or revise a rejected interpretation.

This paper makes three methodological contributions:
\begin{itemize}
  \item We formulate conversational memory retrieval around conjunctive
    plans and source-linked witnesses, separating the interpretation
    of a question from the execution of its relational requirements.
  \item We develop a controller that combines evidence-grounded planning,
    reference-relative ranking, consistent-binding search, and
    item-wise verification with gap-driven replanning.
  \item We describe a compact context construction procedure that uses
    witness support and query relevance to select dated propositions
    and complementary passage summaries for the final reader.
\end{itemize}
% Add empirical contributions only after validating the completed runs,
% comparison protocol, uncertainty estimates, and reported configuration.

\Cref{sec:background} reviews the relevant foundations and prior work.
\Cref{sec:problem,sec:method} define the task and present the method.
The evaluation is organized in \cref{sec:setup,sec:results};
\cref{sec:discussion} discusses limitations.
